\documentclass[11pt]{article}
\usepackage[margin=1in]{geometry}
\usepackage{amsmath,amssymb}
\usepackage{booktabs}
\usepackage{hyperref}
\usepackage{siunitx}

\title{Replication Report: \emph{Skyrmion collapse}\\
\large A.\,D.\ Verga, arXiv:1409.0256v2 (2014)}
\author{REPLICATE-PROJECT / TEXTURES-100 --- textures-polar-verga2014}
\date{Verdict: \textbf{PARTIAL} \quad (Coverage 6/10, Agreement 8/10)}

\begin{document}
\maketitle

\begin{abstract}
Verga's paper is a \emph{dynamics/mechanism} study: a Belavin--Polyakov (BP)
skyrmion driven to a \emph{topological collapse} by a spin-polarized current,
modelled by a Schr\"odinger equation for the itinerant electrons coupled to a
Landau--Lifshitz (LL) equation for the localized spins on a square lattice.
There is \textbf{no Dzyaloshinskii--Moriya interaction (DMI)} in the model ---
the ingredients are Heisenberg exchange, spin-transfer torque (STT), and a
polarization (b-)field. We replicate the paper's \emph{static + scaling core}
from the equations, from scratch, with no author code: (a) the discrete BP
exchange energy converges to the exact continuum value $E_{\mathrm{xc}}=4\pi J$;
(b) the Berg--L\"uscher lattice topological charge is $Q\approx-1$ with the
correct sign; (c) the self-similar collapse exponents $\alpha=1,\ \beta=1/2$
follow from dimensional balance of the paper's ansatz (Eq.~17). We did
\emph{not} build the full coupled Schr\"odinger$+$LL time integrator, the
collapse times $t^*$, or the electron b-field vortex nucleation --- these are
scoped out under the wall-clock budget and listed as honest gaps.
\end{abstract}

\section{The paper and its central claims}

The equilibrium magnetization of nanodots and chiral magnets can carry nontrivial
topology (vortices, skyrmions). Spin-transfer torque from a spin-polarized current
can trigger a \emph{topological change} between states of different topology
without a magnetic field. Verga models this self-consistently: electrons obey a
quantum (Schr\"odinger) dynamics on a square lattice, and the localized texture
$\mathbf S(\mathbf r,t)$, $|\mathbf S|=1$, obeys the Landau--Lifshitz equation
driven by the electron spin density via STT. The magnetic energy is the classical
Heisenberg exchange
\begin{equation}
H_S=\frac{J}{2}\int(\nabla\mathbf S)^2\,d^2r
\;\;\longleftrightarrow\;\;
H=J\!\!\sum_{\langle ij\rangle}\!\bigl(1-\mathbf S_i\!\cdot\!\mathbf S_j\bigr),
\qquad J=0.4 .
\end{equation}
The headline claims we target:
\begin{enumerate}
\item The continuum BP skyrmion $w_0=z/\lambda$ has exchange energy
      $E_{\mathrm{xc}}=4\pi J$ and topological charge $Q=\pm1$, and is
      \emph{scale-invariant} (energy independent of size $\lambda$). (lines 428--429)
\item Without dissipation the coupled system conserves $|\mathbf S|=1$ and $Q(t)$;
      the collapse on the lattice proceeds ``by change of a single spin,'' violating
      $Q$ conservation. (line 179 ff.)
\item The core collapse in the continuum dissipationless limit is a
      \emph{self-similar finite-time singularity}
      $w=(t^*-t)^{-\alpha}f\!\big(r/(t^*-t)^{\beta}\big)$ with
      $\alpha=1,\ \beta=1/2$. (Eq.~17, lines 556--604)
\end{enumerate}

\section{What we built (from scratch, no author code)}

All checks are in \texttt{report/evidence/verga2014\_repl.py} (numpy only), run
with \texttt{/home/stevens/comfyui-env/bin/python}. Outputs in
\texttt{report/evidence/verga2014\_result.json}.

\subsection{(A) Discrete BP exchange energy $\to 4\pi J$}
We place the BP field
$S_z=(\lambda^2-r^2)/(\lambda^2+r^2)$,
$S_x+iS_y=2\lambda(x+iy)/(\lambda^2+r^2)$
on an $L=512$ lattice and evaluate the discrete Heisenberg energy
$H=J\sum_{\langle ij\rangle}(1-\mathbf S_i\!\cdot\!\mathbf S_j)$, the lattice form
of $(J/2)(\nabla\mathbf S)^2$. The continuum benchmark is $4\pi J=5.0265$.

\begin{table}[h]
\centering
\begin{tabular}{rccc}
\toprule
$\lambda$ & $E$ & $E/(4\pi J)$ & $Q$ (Berg--L\"uscher) \\
\midrule
4  & 4.9738 & 0.9895 & $-0.99980$ \\
8  & 5.0095 & 0.9966 & $-0.99920$ \\
16 & 5.0072 & 0.9962 & $-0.99680$ \\
32 & 4.9622 & 0.9872 & $-0.98733$ \\
64 & 4.7817 & 0.9513 & $-0.95124$ \\
\bottomrule
\end{tabular}
\caption{Discrete BP exchange energy vs.\ the continuum $4\pi J$, and lattice
topological charge, on $L=512$. For $\lambda=8$--$16$ the energy matches the
continuum value to $<0.5\%$. The falloff at $\lambda=64$ is the finite-box edge
effect ($\lambda\!\to\!L/8$), not a physics discrepancy.}
\end{table}

\subsection{(B) Lattice topological charge $Q=\pm1$}
Using the Berg--L\"uscher signed solid angle of each spin triangle (gauge-exact,
\emph{not} a naive finite-difference of
$\mathbf S\!\cdot\!(\partial_x\mathbf S\times\partial_y\mathbf S)$), an
$L=256,\ \lambda=20$ field gives $Q=-0.980$ for the core-$-z$ configuration and
$Q=+0.980$ for the reversed field. The sign convention matches the paper
(core $-z\Rightarrow Q=-1$ for the initial state that collapses).

\subsection{(C) Lattice breaks scale invariance $\Rightarrow$ collapse drive}
On a perfect continuum $E(\lambda)=4\pi J$ for all $\lambda$. On the lattice the
energy depends on $\lambda$: it peaks near $\lambda\simeq6$
($E/4\pi J=0.988$) and \emph{decreases} as the core shrinks below a few lattice
sites toward the ferromagnetic ($E=0$) state. There is \textbf{no energy barrier}
protecting the small skyrmion under pure exchange, so a single-spin flip
($\lambda\!\to\!0$) is downhill --- exactly the paper's ``collapse by change of a
single spin,'' regularized only by the lattice cutoff $a$.

\subsection{(D) Self-similar exponents $\alpha=1,\ \beta=1/2$}
The paper's ansatz $w=(t^*-t)^{-\alpha}f(r/(t^*-t)^{\beta})$ inserted into the
driven LL equation gives three scaling powers: time-derivative
$(t^*-t)^{-\alpha-1}$, exchange $\nabla^2w\to(t^*-t)^{-\alpha-2\beta}$, STT
nonlinearity $w^2\to(t^*-t)^{-2\alpha}$. Balancing time-derivative against
exchange and against STT yields the $2\times2$ linear system
\begin{equation}
-\alpha-1=-\alpha-2\beta,\qquad -\alpha-1=-2\alpha
\;\;\Longrightarrow\;\;
\alpha=1,\quad\beta=\tfrac12,
\end{equation}
solved numerically ($\alpha_{\text{computed}}=1.000,\
\beta_{\text{computed}}=0.500$), matching the paper's Eq.~(17). This reproduces
the headline scaling with zero dynamics simulated. The consistent perturbative
size law $\lambda(t)=\lambda_0/\sqrt{1+(s_0t)^2}$ (line 551) has the same
$\beta=1/2$ near collapse.

\section{Comparison to the paper}

\begin{table}[h]
\centering
\begin{tabular}{lccc}
\toprule
Quantity & Paper & This work & Status \\
\midrule
Exchange energy $E_{\mathrm{xc}}$          & $4\pi J=5.0265$ & $5.009$ ($\lambda=8$) & $\checkmark$ $<0.5\%$ \\
Topological charge $Q$                      & $\pm1$          & $\mp0.980$            & $\checkmark$ correct sign \\
Scale-invariance breaking $\to$ collapse    & no barrier      & no barrier            & $\checkmark$ mechanism \\
Self-similar $\alpha$                        & $1$             & $1.000$               & $\checkmark$ \\
Self-similar $\beta$                         & $1/2$           & $0.500$               & $\checkmark$ \\
Coupled Schr\"odinger$+$LL time series       & Figs.~2--3      & \emph{not built}      & gap \\
Collapse time $t^*\sim\lambda/(s_0a)$        & $\sim5900\,t_0$ & \emph{not built}      & gap \\
Electron b-field vortex nucleation           & Sec.~III        & \emph{not built}      & gap \\
\bottomrule
\end{tabular}
\caption{Per-claim comparison. Static energetics, the topological invariant, the
lattice collapse mechanism, and the self-similar exponents all reproduce; the
full coupled time integrator and its derived observables are scoped out.}
\end{table}

\section{Critique and physical assessment}

The static and scaling physics is robust and reproduces cleanly, which is the
right target for a mechanism paper under a wall-clock budget: the checkable
physics lives in the closed-form energy ($4\pi J$), the topological invariant
(Berg--L\"uscher $Q$), and the exponent algebra --- none of which require the
expensive coupled solver. Two points of note:
\begin{itemize}
\item \textbf{No DMI.} The task brief used ``DMI'' generically; the paper contains
      \emph{no} Dzyaloshinskii--Moriya term. Its skyrmion is the pure-exchange BP
      soliton, stabilized dynamically by the current, not energetically by DMI.
      We replicated what the paper does and did not bolt on a DMI term.
\item \textbf{Finite-box falloff is not disagreement.} The $E/(4\pi J)$ drop at
      $\lambda=64$ on $L=512$ is the skyrmion tail reaching the open boundary,
      not a failure of the $4\pi J$ result; the $\lambda=8$--$16$ window is the
      clean comparison.
\end{itemize}

\section{Verdict}

\textbf{PARTIAL} --- Coverage 6/10, Agreement 8/10. Replicated: BP exchange energy
$4\pi J$, quantized topological charge $Q=\pm1$ with correct sign, lattice
scale-invariance breaking driving barrier-free collapse, and the self-similar
exponents $\alpha=1,\ \beta=1/2$. Not replicated (honest gaps): the coupled
Schr\"odinger$+$Landau--Lifshitz time integration with itinerant electrons, the
collapse times $t^*$, and the electron b-field topological-vortex nucleation. See
\texttt{report/failure\_analysis.md} and \texttt{report/open\_questions.json}.

\end{document}
